% !TEX program = pdflatex
\documentclass[10pt,aspectratio=169]{beamer}

\usepackage[utf8]{inputenc}
\usepackage[T1]{fontenc}
\usepackage[english]{babel}
\usepackage{lmodern}
\usepackage{microtype}
\usepackage{amsmath,amssymb,bm,mathtools}
\usepackage{booktabs,array}

\usetheme[numbering=fraction,progressbar=frametitle]{metropolis}
\setbeamertemplate{frame numbering}[fraction]
\setbeamertemplate{navigation symbols}{}

\definecolor{CourseBlue}{HTML}{1A3C68}
\definecolor{CourseGreen}{HTML}{216E39}
\definecolor{CourseGold}{HTML}{B8860B}
\definecolor{CourseRed}{HTML}{9A2C2C}
\definecolor{SoftBlue}{HTML}{F2F6FA}
\definecolor{SoftGreen}{HTML}{F1F8F3}
\definecolor{SoftRed}{HTML}{FFF3F3}

\setbeamercolor{normal text}{fg=black,bg=white}
\setbeamercolor{structure}{fg=CourseBlue}
\setbeamercolor{title}{fg=CourseBlue}
\setbeamerfont{title}{size=\Large,series=\bfseries}
\setbeamercolor{frametitle}{bg=CourseBlue,fg=white}
\setbeamerfont{frametitle}{series=\bfseries,size=\large}
\setbeamertemplate{frametitle}{
  \nointerlineskip
  \begin{beamercolorbox}[wd=\paperwidth,ht=3.2ex,dp=1.4ex,leftskip=1em,rightskip=1em]{frametitle}
    \usebeamerfont{frametitle}\insertframetitle
  \end{beamercolorbox}
}
\setbeamercolor{block title}{bg=CourseBlue,fg=white}
\setbeamercolor{block body}{bg=SoftBlue,fg=black}
\setbeamercolor{block title example}{bg=CourseGreen,fg=white}
\setbeamercolor{block body example}{bg=SoftGreen,fg=black}
\setbeamercolor{block title alerted}{bg=CourseRed,fg=white}
\setbeamercolor{block body alerted}{bg=SoftRed,fg=black}
\setbeamertemplate{blocks}[rounded][shadow=false]

\newenvironment{modelbox}{\begin{beamercolorbox}[rounded=true,sep=0.7em,wd=\textwidth]{block body}
  \textbf{\color{CourseBlue}MODEL}\quad}{\end{beamercolorbox}}
\newenvironment{marketbox}{\begin{beamercolorbox}[rounded=true,sep=0.7em,wd=\textwidth]{block body example}
  \textbf{\color{CourseGreen}NO-ARBITRAGE}\quad}{\end{beamercolorbox}}
\newenvironment{simplificationbox}{\begin{beamercolorbox}[rounded=true,sep=0.7em,wd=\textwidth]{block body alerted}
  \textbf{\color{CourseRed}WATCH OUT}\quad}{\end{beamercolorbox}}
\newenvironment{pausebox}{\begin{beamercolorbox}[rounded=true,sep=0.8em,wd=\textwidth]{block body}
  \textbf{\color{CourseGold}YOUR TURN}\quad}{\end{beamercolorbox}}

\newcommand{\E}{\mathbb E}
\newcommand{\Pp}{\mathbb P}
\newcommand{\Q}{\mathbb Q}
\newcommand{\F}{\mathcal F}
\newcommand{\dd}{\,\mathrm d}
\newcommand{\Ncdf}{\Phi}

\AtBeginSection[]{
  \begin{frame}[plain]\vfill
    {\Large\bfseries\color{CourseBlue}\insertsectionhead\par}
    \medskip{\color{CourseBlue}\rule{0.92\textwidth}{0.8pt}}\vfill
  \end{frame}
}

\title{Pricing Toolbox: A Guided Refresher}
\subtitle{It\^o calculus, Girsanov, path dependence, and early exercise}
\author{}
\date{Lecture 14}

\begin{document}

\begin{frame}[plain]
  \vfill
  {\LARGE\bfseries\color{CourseBlue}Pricing Toolbox: A Guided Refresher\par}
  \medskip{\large It\^o calculus, Girsanov, path dependence, and early exercise\par}
  \medskip{\color{CourseGold}\rule{0.94\textwidth}{0.8pt}}
  \medskip{Lecture 14\par}\vfill
\end{frame}

\begin{frame}{Four tools solve four pricing questions}
\begin{block}{The question}
How do the core tools of continuous-time finance work together when a payoff
depends on a path, a barrier, or an exercise decision?
\end{block}
\begin{enumerate}
  \item Use It\^o's formula to handle nonlinear functions of Brownian motion.
  \item Use Girsanov to move a drifted hitting problem to Brownian motion.
  \item Augment the state for an Asian option and price it by PDE or simulation.
  \item Use backward induction to decide whether an American option is exercised.
\end{enumerate}
\begin{marketbox}
Each technique changes the representation of a problem while preserving the
economic object that must be priced.
\end{marketbox}
\end{frame}

\section{It\^o calculus turns stochastic integrals into identities}

\begin{frame}{The quadratic-variation term is the essential correction}
For \(f(x)=e^x\), It\^o's formula gives
\[
\dd e^{B_t}=e^{B_t}\dd B_t+\frac12e^{B_t}\dd t.
\]
Integrating from \(0\) to \(t\),
\[
\boxed{
\int_0^t e^{B_s}\dd B_s
=
e^{B_t}-1-\frac12\int_0^t e^{B_s}\dd s.}
\]
\begin{block}{Why ordinary calculus fails}
The term \(\tfrac12 f''(B_t)\dd t\) comes from
\((\dd B_t)^2=\dd t\). Omitting it produces the wrong stochastic integral.
\end{block}
\end{frame}

\begin{frame}{A quick It\^o check}
\begin{pausebox}
Suppose someone claims
\[
\int_0^t e^{B_s}\dd B_s=e^{B_t}-1.
\]
What term is missing, and why?
\end{pausebox}
\pause
\begin{exampleblock}{Answer}
The correction is
\[
-\frac12\int_0^t e^{B_s}\dd s.
\]
It comes from the second derivative of \(e^x\) and Brownian quadratic
variation.
\end{exampleblock}
\end{frame}

\begin{frame}{The stochastic exponential is a martingale after correction}
For deterministic square-integrable \(f\), define
\[
\mathcal E_t=
\exp\left(
\int_0^t f(s)\dd B_s-\frac12\int_0^t f(s)^2\dd s
\right).
\]
Then
\[
\dd\mathcal E_t=\mathcal E_t f(t)\dd B_t,
\]
so \(\mathcal E\) is a martingale under standard integrability conditions.
\begin{marketbox}
The compensating term
\(\tfrac12\int f^2\) is exactly what makes an exponential Gaussian
quantity have mean one. It will become the density process in Girsanov.
\end{marketbox}
\end{frame}

\section{Girsanov removes drift by changing measure}

\begin{frame}{A change of measure turns drifted Brownian motion into Brownian motion}
Let
\[
X_t=\mu t+\sigma B_t,\qquad \theta=\frac{\mu}{\sigma},
\]
under \(\Pp\). Define
\[
Z_t=\exp\left(-\theta B_t-\frac12\theta^2t\right),
\qquad
\left.\frac{\dd\Q}{\dd\Pp}\right|_{\F_t}=Z_t.
\]
Girsanov's theorem says
\[
\widetilde B_t=B_t+\theta t
\]
is a Brownian motion under \(\Q\). Therefore
\[
\boxed{X_t=\sigma\widetilde B_t\quad\text{under }\Q.}
\]
\end{frame}

\begin{frame}{The measure change simplifies a hitting-time question}
Let
\[
\tau_a=\inf\{t\geq0:X_t=a\},\qquad a>0.
\]
Under \(\Q\), \(\tau_a\) is the first time a standard Brownian motion reaches
\(a/\sigma\). Its density is
\[
f_{\tau_a}^{\Q}(t)
=
\frac{a/\sigma}{\sqrt{2\pi t^3}}
\exp\left(-\frac{(a/\sigma)^2}{2t}\right).
\]
\begin{block}{Returning to the original measure}
Use
\[
\Pp(A)=\E^\Q[Z_t^{-1}\mathbf1_A],
\]
then evaluate the density weight on the event \(\{\tau_a=t\}\).
\end{block}
\end{frame}

\begin{frame}{A drifted hitting time has an inverse-Gaussian density}
The change-of-measure calculation yields
\[
\boxed{
f_{\tau_a}^{\Pp}(t)=
\frac{a}{\sigma\sqrt{2\pi t^3}}
\exp\left(-\frac{(a-\mu t)^2}{2\sigma^2t}\right),\qquad t>0.}
\]
\begin{pausebox}
When \(\mu>0\), does the process tend to hit a positive level sooner or later
than a driftless Brownian motion?
\end{pausebox}
\pause
\begin{exampleblock}{Answer}
Sooner. Positive drift shifts the density toward earlier hitting times. The
formula makes this visible through the term \((a-\mu t)^2\).
\end{exampleblock}
\end{frame}

\section{Martingale pricing with path-dependent payoffs}

\begin{frame}{An Asian payoff needs a memory variable}
An arithmetic-average Asian call pays
\[
\left(\frac1T\int_0^T S_u\dd u-K\right)^+.
\]
Define the running integral
\[
I_t=\int_0^t S_u\dd u.
\]
At time \(t\),
\[
\int_0^T S_u\dd u=I_t+\int_t^T S_u\dd u.
\]
\begin{marketbox}
The stock price alone is not enough. The minimal Markov state is
\(\boxed{(t,S_t,I_t)}\).
\end{marketbox}
\end{frame}

\begin{frame}{State augmentation restores a pricing PDE}
Under Black--Scholes,
\[
\dd S_t=rS_t\dd t+\sigma S_t\dd B_t,\qquad
\dd I_t=S_t\dd t.
\]
For \(v(t,s,i)\), the pricing PDE is
\[
\boxed{
v_t+rs\,v_s+s\,v_i+\tfrac12\sigma^2s^2v_{ss}-rv=0,}
\]
with terminal condition
\[
v(T,s,i)=\left(\frac{i}{T}-K\right)^+.
\]
\begin{block}{What changed from a vanilla option}
The new term \(s\,v_i\) reflects the deterministic accumulation of the
running integral; there is no diffusion term in the \(i\)-direction.
\end{block}
\end{frame}

\begin{frame}{Monte Carlo prices the same Asian claim path by path}
On a grid \(t_k=k\Delta t\), simulate risk-neutral stock paths and form
\[
\bar S_T^{(m)}=\frac1T\sum_{k=0}^{n-1}S_{t_k}^{(m)}\Delta t.
\]
The price estimator is
\[
\boxed{V_0\approx e^{-rT}\frac1M\sum_{m=1}^M
\left(\bar S_T^{(m)}-K\right)^+.}
\]
\begin{pausebox}
Where should the deterministic discount factor appear?
\end{pausebox}
\pause
\begin{exampleblock}{Answer}
Once outside the average. The simulated quantity is the risk-neutral payoff.
\end{exampleblock}
\end{frame}

\begin{frame}{A geometric Asian option can reduce Monte Carlo noise}
The geometric average
\[
G=\exp\left(\frac1T\int_0^T\log S_u\dd u\right)
\]
is tractable in Black--Scholes because \(\log S\) is Gaussian and its time
integral is Gaussian.
\begin{columns}[T,onlytextwidth]
\column{0.48\textwidth}
\begin{block}{Arithmetic Asian}
The desired payoff uses the arithmetic average; it generally has no
comparable closed form.
\end{block}
\column{0.48\textwidth}
\begin{exampleblock}{Control variate}
Simulate both payoffs and use the known geometric-Asian value to reduce
variance in the arithmetic-Asian estimator.
\end{exampleblock}
\end{columns}
\end{frame}

\section{Barrier events combine a path maximum and Girsanov}

\begin{frame}{A one-touch option is a hitting probability}
Under \(\Q\), let the stock pay dividend yield \(q\):
\[
\frac{\dd S_t}{S_t}=(r-q)\dd t+\sigma\dd B_t.
\]
A one-touch pays one unit at \(T\) if it ever reaches \(H>S_0\):
\[
\Pi=\mathbf1_{\{\sup_{0\leq t\leq T}S_t\geq H\}}.
\]
\begin{block}{Log transform}
With \(X_t=\log(S_t/S_0)=\mu t+\sigma B_t\), where
\(\mu=(r-q)-\tfrac12\sigma^2\), the barrier becomes the level
\(a=\log(H/S_0)\). Price the resulting hitting probability and discount it.
\end{block}
\end{frame}

\begin{frame}{Reflection plus Girsanov gives the one-touch probability}
For \(Y_t=\mu t+\sigma B_t\),
\[
\boxed{
\Q\left(\sup_{0\leq t\leq T}Y_t\geq a\right)
=
\Phi\left(\frac{\mu T-a}{\sigma\sqrt T}\right)
+e^{2\mu a/\sigma^2}
\Phi\left(\frac{-\mu T-a}{\sigma\sqrt T}\right).}
\]
\begin{columns}[T,onlytextwidth]
\column{0.48\textwidth}
\begin{block}{When \(\mu=0\)}
Reflection gives twice the upper tail of \(B_T\).
\end{block}
\column{0.48\textwidth}
\begin{exampleblock}{When \(\mu\neq0\)}
Girsanov removes drift, then the density change restores the original
probability weights.
\end{exampleblock}
\end{columns}
\end{frame}

\begin{frame}{Worked one-touch: check the scale of a barrier probability}
Let \(S_0=100\), \(H=120\), \(r=5\%\), \(q=0\), \(\sigma=20\%\), \(T=1\).
Then
\[
\mu=0.03,\qquad a=\log(1.2)=0.1823,
\]
and the hitting probability is approximately
\[
\Q\left(\sup_{t\leq1}S_t\geq120\right)\approx0.413.
\]
Therefore
\[
\boxed{V_0\approx e^{-0.05}(0.413)=0.393.}
\]
\begin{marketbox}
A terminal digital and a one-touch are not the same contract: the barrier
may be hit and then the stock may fall before maturity.
\end{marketbox}
\end{frame}

\begin{frame}{Discrete monitoring understates continuous barrier hitting}
A simulation on dates \(t_0,\ldots,t_N\) replaces
\[
\mathbf1_{\{\sup_{0\leq t\leq T}S_t\geq H\}}
\]
by
\[
\mathbf1_{\{\max_{0\leq k\leq N}S_{t_k}\geq H\}}.
\]
\begin{pausebox}
For an up-and-in one-touch, is this discretely monitored estimator biased
high or low relative to continuous monitoring?
\end{pausebox}
\pause
\begin{exampleblock}{Answer}
Low. A path can cross the barrier between two monitoring dates and return
below it before the next observation. Refining the grid reduces this bias.
\end{exampleblock}
\end{frame}

\section{American options turn pricing into a decision}

\begin{frame}{American pricing is backward induction with an exercise test}
In a binomial model with gross return \(R=1+r\) and risk-neutral probability
\(p\),
\[
\boxed{
V_t=\max\left\{g(S_t),\frac{pV_{t+1}^{u}+(1-p)V_{t+1}^{d}}{R}\right\}.}
\]
\begin{columns}[T,onlytextwidth]
\column{0.48\textwidth}
\begin{block}{Exercise value}
\[
g(S_t).
\]
Take the payoff and stop.
\end{block}
\column{0.48\textwidth}
\begin{exampleblock}{Continuation value}
\[
R^{-1}\E^\Q[V_{t+1}\mid\F_t].
\]
Keep all future exercise choices.
\end{exampleblock}
\end{columns}
\end{frame}

\begin{frame}{A three-step American put has a genuine exercise region}
Consider
\[
S_0=100,\quad u=1.1,\quad d=0.9,\quad R=1.02,\quad K=100.
\]
The risk-neutral up probability is
\[
p=\frac{R-d}{u-d}=0.6.
\]
At maturity, the stock values and put payoffs are
\[
\begin{array}{c|cccc}
S_3 & 133.1 & 108.9 & 89.1 & 72.9\\
\midrule
(K-S_3)^+ & 0 & 0 & 10.9 & 27.1
\end{array}
\]
\begin{marketbox}
At \(S_2=81\), \(19>
\bigl[0.6(10.9)+0.4(27.1)\bigr]/1.02=17.04\), so exercise.
\end{marketbox}
\end{frame}

\begin{frame}{The early-exercise premium comes from node-by-node choices}
Working backward gives:
\[
\begin{array}{c|c|c|c}
\toprule
\text{Node} & \text{Exercise} & \text{Continuation} & \text{Action}\\
\midrule
t=2,\ S=81 & 19.00 & 17.04 & exercise\\
t=1,\ S=90 & 10.00 & 9.97 & exercise\\
t=0,\ S=100 & 0.00 & 4.91 & continue\\
\bottomrule
\end{array}
\]
\[
V_0^{\mathrm{Eur}}\approx4.59,\qquad V_0^{\mathrm{Am}}\approx4.91.
\]
\begin{pausebox}
What is the early-exercise premium?
\end{pausebox}
\pause
\begin{exampleblock}{Answer}
\(4.91-4.59=0.32\): the value of exercising at low-stock nodes.
\end{exampleblock}
\end{frame}

\section{Numerical habits complete the toolkit}

\begin{frame}{Euler simulation is simple, but its approximation must be named}
For
\[
\dd X_t=b(X_t,t)\dd t+a(X_t,t)\dd W_t,
\]
the Euler scheme on a grid is
\[
\boxed{
X_{n+1}=X_n+b(X_n,t_n)\Delta t
+a(X_n,t_n)\sqrt{\Delta t}\,Z_{n+1},}
\]
where \(Z_{n+1}\sim N(0,1)\).
\begin{columns}[T,onlytextwidth]
\column{0.48\textwidth}
\begin{block}{What is approximated}
The continuous diffusion and any path functional evaluated between grid
dates.
\end{block}
\column{0.48\textwidth}
\begin{alertblock}{What must be checked}
Convergence in \(\Delta t\), sampling error in \(M\), positivity or boundary
behavior, and bias for extrema or stopping rules.
\end{alertblock}
\end{columns}
\end{frame}

\begin{frame}{Before trusting a numerical price, identify its weakest link}
\begin{center}
\begin{tabular}{p{0.31\textwidth}|p{0.49\textwidth}}
\toprule
Problem & First diagnostic\\
\midrule
Asian Monte Carlo & refine the integration grid and estimate sampling error\\
Barrier Monte Carlo & test monitoring bias and barrier correction\\
American tree & refine time steps and map the exercise region\\
PDE or local-vol grid & test boundary placement and mesh convergence\\
\bottomrule
\end{tabular}
\end{center}
\begin{marketbox}
A numerical result is credible only when it survives a refinement or a
no-arbitrage check appropriate to the contract.
\end{marketbox}
\end{frame}

\section{Takeaways}

\begin{frame}{The same pricing logic adapts to the contract}
\begin{center}
\begin{tabular}{p{0.25\textwidth}|p{0.30\textwidth}|p{0.27\textwidth}}
\toprule
Contract feature & Mathematical tool & Extra state or decision\\
\midrule
Nonlinear Brownian function & It\^o formula & quadratic variation\\
Drifted barrier event & Girsanov and reflection & hitting probability\\
Asian average & state augmentation & running integral \(I_t\)\\
American exercise & dynamic programming & stop or continue\\
\bottomrule
\end{tabular}
\end{center}
\begin{marketbox}
The payoff determines which information must be carried forward: a state
variable, a stopping rule, or a path-monitoring convention.
\end{marketbox}
\end{frame}

\begin{frame}{Final check: choose the representation before calculating}
\begin{pausebox}
A claim pays one dollar at maturity if a stock ever crosses an upper barrier.
What are the three modelling choices you should state before presenting a
numerical price?
\end{pausebox}
\pause
\begin{exampleblock}{Answer}
State the risk-neutral dynamics, whether monitoring is continuous or
discrete, and the numerical approximation used for the hitting event. These
choices determine the probability being priced.
\end{exampleblock}
\vfill
\begin{center}
{\large\color{CourseBlue}\bfseries Good pricing begins by identifying the right state, measure, and event.}
\end{center}
\end{frame}

\end{document}
